% Shared preamble for all spec documents.
% Include from each spec: % Shared preamble for all spec documents.
% Include from each spec: % Shared preamble for all spec documents.
% Include from each spec: \input{../preamble.tex} (adjust depth).

\documentclass[11pt,a4paper]{article}

\usepackage[utf8]{inputenc}
\usepackage[T1]{fontenc}
\usepackage{lmodern}
\usepackage[a4paper,margin=2.3cm]{geometry}
\usepackage{parskip}
\usepackage{microtype}
\usepackage[dvipsnames]{xcolor}
\usepackage{enumitem}
\usepackage{listings}
\usepackage{tcolorbox}
\usepackage{titlesec}
\usepackage{fancyhdr}
\usepackage{array}
\usepackage{longtable}
\usepackage{booktabs}
\usepackage{amsmath}
\usepackage{amssymb}
\usepackage{hyperref}
\usepackage{cleveref}

\tcbuselibrary{skins,breakable}

\definecolor{specBlue}{RGB}{30,64,132}
\definecolor{specGray}{RGB}{90,90,90}
\definecolor{specAccent}{RGB}{198,40,40}
\definecolor{specGreen}{RGB}{30,110,60}

\hypersetup{
  colorlinks=true,
  linkcolor=specBlue,
  urlcolor=specBlue,
  citecolor=specBlue,
  pdfborder={0 0 0}
}

\titleformat{\section}{\Large\bfseries\color{specBlue}}{\thesection}{0.75em}{}
\titleformat{\subsection}{\large\bfseries\color{specBlue!80}}{\thesubsection}{0.6em}{}

\makeatletter
\newcommand{\specID}[1]{\def\@specID{#1}}
\newcommand{\specTitle}[1]{\def\@specTitle{#1}}
\newcommand{\specStatus}[1]{\def\@specStatus{#1}}
\newcommand{\specVersion}[1]{\def\@specVersion{#1}}
\newcommand{\specOwner}[1]{\def\@specOwner{#1}}
\newcommand{\specTraces}[1]{\def\@specTraces{#1}}

\specID{SPEC-XXX}
\specTitle{Untitled Spec}
\specStatus{DRAFT}
\specVersion{0.1}
\specOwner{Jeremie Nunez}
\specTraces{}

\newcommand{\makespecheader}{%
  \begin{center}
    {\LARGE\bfseries\color{specBlue}\@specTitle}\\[0.4em]
    {\small\color{specGray}\texttt{\@specID} \quad v\@specVersion \quad \textsc{\@specStatus} \quad \@specOwner}
  \end{center}
  \vspace{0.4em}
  \hrule
  \vspace{0.8em}
}

\pagestyle{fancy}
\fancyhf{}
\fancyhead[L]{\small\color{specGray}\@specID}
\fancyhead[R]{\small\color{specGray}\@specTitle}
\fancyfoot[C]{\small\color{specGray}\thepage}
\renewcommand{\headrulewidth}{0pt}
\makeatother

\newtcolorbox{invariant}{
  colback=specBlue!5, colframe=specBlue!75,
  title=Invariant, fonttitle=\bfseries, breakable,
  boxrule=0.5pt, arc=2pt
}

\newtcolorbox{scenario}[1]{
  colback=white, colframe=specBlue!50,
  title={Scenario: #1}, fonttitle=\bfseries, breakable,
  boxrule=0.5pt, arc=2pt
}

\newtcolorbox{ac}[1]{
  colback=specAccent!5, colframe=specAccent!60,
  title={Acceptance Criterion #1}, fonttitle=\bfseries, breakable,
  boxrule=0.5pt, arc=2pt
}

\newtcolorbox{nongoal}{
  colback=specGray!8, colframe=specGray!60,
  title=Non-Goal, fonttitle=\bfseries, breakable,
  boxrule=0.5pt, arc=2pt
}

\lstdefinestyle{codestyle}{
  basicstyle=\ttfamily\small,
  breaklines=true,
  showstringspaces=false,
  frame=single,
  framesep=4pt,
  rulecolor=\color{specBlue!30},
  backgroundcolor=\color{specBlue!2},
  xleftmargin=0.5em,
  xrightmargin=0.5em,
  keywordstyle=\color{specBlue}\bfseries,
  commentstyle=\color{specGray}\itshape,
  stringstyle=\color{specGreen}
}
\lstset{
  inputencoding=utf8,
  extendedchars=true,
  literate=
    {—}{{--}}1
    {–}{{-}}1
    {…}{{...}}1
    {’}{{'}}1
    {‘}{{'}}1
    {“}{{``}}1
    {”}{{''}}1
    {é}{{\'e}}1
    {è}{{\`e}}1
    {ê}{{\^e}}1
    {à}{{\`a}}1
    {â}{{\^a}}1
    {ç}{{\c{c}}}1
    {ô}{{\^o}}1
    {→}{{$\to$}}1
    {←}{{$\leftarrow$}}1
    {×}{{$\times$}}1
    {≥}{{$\geq$}}1
    {≤}{{$\leq$}}1
    {≠}{{$\neq$}}1
    {∈}{{$\in$}}1
    {⌈}{{$\lceil$}}1
    {⌉}{{$\rceil$}}1
    {∞}{{$\infty$}}1
    {α}{{$\alpha$}}1
    {β}{{$\beta$}}1
}
\lstdefinelanguage{TypeScript}{
  keywords={const,let,var,function,return,if,else,for,while,class,interface,type,extends,implements,import,export,from,as,async,await,new,this,true,false,null,undefined,void,any,unknown,never,string,number,boolean,readonly,private,public,protected,enum},
  sensitive=true,
  comment=[l]{//},
  morecomment=[s]{/*}{*/},
  morestring=[b]",
  morestring=[b]',
  morestring=[b]`
}
\lstset{style=codestyle}

\newcommand{\Given}{\textbf{\color{specBlue}Given} }
\newcommand{\When}{\textbf{\color{specBlue}When} }
\newcommand{\Then}{\textbf{\color{specBlue}Then} }
\providecommand{\And}{}
\renewcommand{\And}{\textbf{\color{specBlue}And} }
 (adjust depth).

\documentclass[11pt,a4paper]{article}

\usepackage[utf8]{inputenc}
\usepackage[T1]{fontenc}
\usepackage{lmodern}
\usepackage[a4paper,margin=2.3cm]{geometry}
\usepackage{parskip}
\usepackage{microtype}
\usepackage[dvipsnames]{xcolor}
\usepackage{enumitem}
\usepackage{listings}
\usepackage{tcolorbox}
\usepackage{titlesec}
\usepackage{fancyhdr}
\usepackage{array}
\usepackage{longtable}
\usepackage{booktabs}
\usepackage{amsmath}
\usepackage{amssymb}
\usepackage{hyperref}
\usepackage{cleveref}

\tcbuselibrary{skins,breakable}

\definecolor{specBlue}{RGB}{30,64,132}
\definecolor{specGray}{RGB}{90,90,90}
\definecolor{specAccent}{RGB}{198,40,40}
\definecolor{specGreen}{RGB}{30,110,60}

\hypersetup{
  colorlinks=true,
  linkcolor=specBlue,
  urlcolor=specBlue,
  citecolor=specBlue,
  pdfborder={0 0 0}
}

\titleformat{\section}{\Large\bfseries\color{specBlue}}{\thesection}{0.75em}{}
\titleformat{\subsection}{\large\bfseries\color{specBlue!80}}{\thesubsection}{0.6em}{}

\makeatletter
\newcommand{\specID}[1]{\def\@specID{#1}}
\newcommand{\specTitle}[1]{\def\@specTitle{#1}}
\newcommand{\specStatus}[1]{\def\@specStatus{#1}}
\newcommand{\specVersion}[1]{\def\@specVersion{#1}}
\newcommand{\specOwner}[1]{\def\@specOwner{#1}}
\newcommand{\specTraces}[1]{\def\@specTraces{#1}}

\specID{SPEC-XXX}
\specTitle{Untitled Spec}
\specStatus{DRAFT}
\specVersion{0.1}
\specOwner{Jeremie Nunez}
\specTraces{}

\newcommand{\makespecheader}{%
  \begin{center}
    {\LARGE\bfseries\color{specBlue}\@specTitle}\\[0.4em]
    {\small\color{specGray}\texttt{\@specID} \quad v\@specVersion \quad \textsc{\@specStatus} \quad \@specOwner}
  \end{center}
  \vspace{0.4em}
  \hrule
  \vspace{0.8em}
}

\pagestyle{fancy}
\fancyhf{}
\fancyhead[L]{\small\color{specGray}\@specID}
\fancyhead[R]{\small\color{specGray}\@specTitle}
\fancyfoot[C]{\small\color{specGray}\thepage}
\renewcommand{\headrulewidth}{0pt}
\makeatother

\newtcolorbox{invariant}{
  colback=specBlue!5, colframe=specBlue!75,
  title=Invariant, fonttitle=\bfseries, breakable,
  boxrule=0.5pt, arc=2pt
}

\newtcolorbox{scenario}[1]{
  colback=white, colframe=specBlue!50,
  title={Scenario: #1}, fonttitle=\bfseries, breakable,
  boxrule=0.5pt, arc=2pt
}

\newtcolorbox{ac}[1]{
  colback=specAccent!5, colframe=specAccent!60,
  title={Acceptance Criterion #1}, fonttitle=\bfseries, breakable,
  boxrule=0.5pt, arc=2pt
}

\newtcolorbox{nongoal}{
  colback=specGray!8, colframe=specGray!60,
  title=Non-Goal, fonttitle=\bfseries, breakable,
  boxrule=0.5pt, arc=2pt
}

\lstdefinestyle{codestyle}{
  basicstyle=\ttfamily\small,
  breaklines=true,
  showstringspaces=false,
  frame=single,
  framesep=4pt,
  rulecolor=\color{specBlue!30},
  backgroundcolor=\color{specBlue!2},
  xleftmargin=0.5em,
  xrightmargin=0.5em,
  keywordstyle=\color{specBlue}\bfseries,
  commentstyle=\color{specGray}\itshape,
  stringstyle=\color{specGreen}
}
\lstset{
  inputencoding=utf8,
  extendedchars=true,
  literate=
    {—}{{--}}1
    {–}{{-}}1
    {…}{{...}}1
    {’}{{'}}1
    {‘}{{'}}1
    {“}{{``}}1
    {”}{{''}}1
    {é}{{\'e}}1
    {è}{{\`e}}1
    {ê}{{\^e}}1
    {à}{{\`a}}1
    {â}{{\^a}}1
    {ç}{{\c{c}}}1
    {ô}{{\^o}}1
    {→}{{$\to$}}1
    {←}{{$\leftarrow$}}1
    {×}{{$\times$}}1
    {≥}{{$\geq$}}1
    {≤}{{$\leq$}}1
    {≠}{{$\neq$}}1
    {∈}{{$\in$}}1
    {⌈}{{$\lceil$}}1
    {⌉}{{$\rceil$}}1
    {∞}{{$\infty$}}1
    {α}{{$\alpha$}}1
    {β}{{$\beta$}}1
}
\lstdefinelanguage{TypeScript}{
  keywords={const,let,var,function,return,if,else,for,while,class,interface,type,extends,implements,import,export,from,as,async,await,new,this,true,false,null,undefined,void,any,unknown,never,string,number,boolean,readonly,private,public,protected,enum},
  sensitive=true,
  comment=[l]{//},
  morecomment=[s]{/*}{*/},
  morestring=[b]",
  morestring=[b]',
  morestring=[b]`
}
\lstset{style=codestyle}

\newcommand{\Given}{\textbf{\color{specBlue}Given} }
\newcommand{\When}{\textbf{\color{specBlue}When} }
\newcommand{\Then}{\textbf{\color{specBlue}Then} }
\providecommand{\And}{}
\renewcommand{\And}{\textbf{\color{specBlue}And} }
 (adjust depth).

\documentclass[11pt,a4paper]{article}

\usepackage[utf8]{inputenc}
\usepackage[T1]{fontenc}
\usepackage{lmodern}
\usepackage[a4paper,margin=2.3cm]{geometry}
\usepackage{parskip}
\usepackage{microtype}
\usepackage[dvipsnames]{xcolor}
\usepackage{enumitem}
\usepackage{listings}
\usepackage{tcolorbox}
\usepackage{titlesec}
\usepackage{fancyhdr}
\usepackage{array}
\usepackage{longtable}
\usepackage{booktabs}
\usepackage{amsmath}
\usepackage{amssymb}
\usepackage{hyperref}
\usepackage{cleveref}

\tcbuselibrary{skins,breakable}

\definecolor{specBlue}{RGB}{30,64,132}
\definecolor{specGray}{RGB}{90,90,90}
\definecolor{specAccent}{RGB}{198,40,40}
\definecolor{specGreen}{RGB}{30,110,60}

\hypersetup{
  colorlinks=true,
  linkcolor=specBlue,
  urlcolor=specBlue,
  citecolor=specBlue,
  pdfborder={0 0 0}
}

\titleformat{\section}{\Large\bfseries\color{specBlue}}{\thesection}{0.75em}{}
\titleformat{\subsection}{\large\bfseries\color{specBlue!80}}{\thesubsection}{0.6em}{}

\makeatletter
\newcommand{\specID}[1]{\def\@specID{#1}}
\newcommand{\specTitle}[1]{\def\@specTitle{#1}}
\newcommand{\specStatus}[1]{\def\@specStatus{#1}}
\newcommand{\specVersion}[1]{\def\@specVersion{#1}}
\newcommand{\specOwner}[1]{\def\@specOwner{#1}}
\newcommand{\specTraces}[1]{\def\@specTraces{#1}}

\specID{SPEC-XXX}
\specTitle{Untitled Spec}
\specStatus{DRAFT}
\specVersion{0.1}
\specOwner{Jeremie Nunez}
\specTraces{}

\newcommand{\makespecheader}{%
  \begin{center}
    {\LARGE\bfseries\color{specBlue}\@specTitle}\\[0.4em]
    {\small\color{specGray}\texttt{\@specID} \quad v\@specVersion \quad \textsc{\@specStatus} \quad \@specOwner}
  \end{center}
  \vspace{0.4em}
  \hrule
  \vspace{0.8em}
}

\pagestyle{fancy}
\fancyhf{}
\fancyhead[L]{\small\color{specGray}\@specID}
\fancyhead[R]{\small\color{specGray}\@specTitle}
\fancyfoot[C]{\small\color{specGray}\thepage}
\renewcommand{\headrulewidth}{0pt}
\makeatother

\newtcolorbox{invariant}{
  colback=specBlue!5, colframe=specBlue!75,
  title=Invariant, fonttitle=\bfseries, breakable,
  boxrule=0.5pt, arc=2pt
}

\newtcolorbox{scenario}[1]{
  colback=white, colframe=specBlue!50,
  title={Scenario: #1}, fonttitle=\bfseries, breakable,
  boxrule=0.5pt, arc=2pt
}

\newtcolorbox{ac}[1]{
  colback=specAccent!5, colframe=specAccent!60,
  title={Acceptance Criterion #1}, fonttitle=\bfseries, breakable,
  boxrule=0.5pt, arc=2pt
}

\newtcolorbox{nongoal}{
  colback=specGray!8, colframe=specGray!60,
  title=Non-Goal, fonttitle=\bfseries, breakable,
  boxrule=0.5pt, arc=2pt
}

\lstdefinestyle{codestyle}{
  basicstyle=\ttfamily\small,
  breaklines=true,
  showstringspaces=false,
  frame=single,
  framesep=4pt,
  rulecolor=\color{specBlue!30},
  backgroundcolor=\color{specBlue!2},
  xleftmargin=0.5em,
  xrightmargin=0.5em,
  keywordstyle=\color{specBlue}\bfseries,
  commentstyle=\color{specGray}\itshape,
  stringstyle=\color{specGreen}
}
\lstset{
  inputencoding=utf8,
  extendedchars=true,
  literate=
    {—}{{--}}1
    {–}{{-}}1
    {…}{{...}}1
    {’}{{'}}1
    {‘}{{'}}1
    {“}{{``}}1
    {”}{{''}}1
    {é}{{\'e}}1
    {è}{{\`e}}1
    {ê}{{\^e}}1
    {à}{{\`a}}1
    {â}{{\^a}}1
    {ç}{{\c{c}}}1
    {ô}{{\^o}}1
    {→}{{$\to$}}1
    {←}{{$\leftarrow$}}1
    {×}{{$\times$}}1
    {≥}{{$\geq$}}1
    {≤}{{$\leq$}}1
    {≠}{{$\neq$}}1
    {∈}{{$\in$}}1
    {⌈}{{$\lceil$}}1
    {⌉}{{$\rceil$}}1
    {∞}{{$\infty$}}1
    {α}{{$\alpha$}}1
    {β}{{$\beta$}}1
}
\lstdefinelanguage{TypeScript}{
  keywords={const,let,var,function,return,if,else,for,while,class,interface,type,extends,implements,import,export,from,as,async,await,new,this,true,false,null,undefined,void,any,unknown,never,string,number,boolean,readonly,private,public,protected,enum},
  sensitive=true,
  comment=[l]{//},
  morecomment=[s]{/*}{*/},
  morestring=[b]",
  morestring=[b]',
  morestring=[b]`
}
\lstset{style=codestyle}

\newcommand{\Given}{\textbf{\color{specBlue}Given} }
\newcommand{\When}{\textbf{\color{specBlue}When} }
\newcommand{\Then}{\textbf{\color{specBlue}Then} }
\providecommand{\And}{}
\renewcommand{\And}{\textbf{\color{specBlue}And} }


\specID{SPEC-SH-002}
\specTitle{AppError hierarchy --- structured error propagation}
\specStatus{DRAFT}
\specVersion{0.1}
\specOwner{Jeremie Nunez}

\begin{document}
\makespecheader

\section{Context}
Errors flowing out of controllers, services, and agent cortices must carry a stable, serialisable contract so that transport layers (HTTP, queue workers, log pipelines) can map them to status codes and structured logs without string-matching messages. The \texttt{AppError} hierarchy in \texttt{packages/shared/src/utils/error.ts} defines a small, closed set of error classes --- \texttt{AppError}, \texttt{ValidationError}, \texttt{NotFoundError}, \texttt{UnauthorizedError}, \texttt{SystemError} --- plus two helpers (\texttt{isAppError}, \texttt{getErrorMessage}) and a log formatter (\texttt{formatError}).

\section{Goals}
\begin{itemize}[leftmargin=1.2em]
  \item Provide a single base class (\texttt{AppError}) carrying a numeric \texttt{statusCode} for transport mapping.
  \item Offer four subclasses with fixed default status codes (400, 404, 401, 500) and distinct \texttt{name} values.
  \item Preserve \texttt{instanceof} relationships across the hierarchy so narrowing works predictably.
  \item Provide a type guard (\texttt{isAppError}) and pure helpers (\texttt{getErrorMessage}, \texttt{formatError}) that accept \texttt{unknown}.
\end{itemize}

\section{Non-Goals}
\begin{nongoal}
This hierarchy does not define transport-level serialisation (JSON body shape, HTTP headers), localisation of messages, or error codes beyond \texttt{statusCode}. It also does not aggregate multiple errors --- each instance describes one failure.
\end{nongoal}

\section{Invariants}
\begin{invariant}
Every subclass \texttt{C} of \texttt{AppError} satisfies: \texttt{(new C(msg)) instanceof AppError \&\& (new C(msg)) instanceof Error}.
\end{invariant}

\begin{invariant}
The \texttt{name} property of each instance equals the class name exactly: \texttt{"AppError"}, \texttt{"ValidationError"}, \texttt{"NotFoundError"}, \texttt{"UnauthorizedError"}, \texttt{"SystemError"}.
\end{invariant}

\begin{invariant}
Default status codes are fixed: \texttt{AppError} = 500, \texttt{ValidationError} = 400, \texttt{NotFoundError} = 404, \texttt{UnauthorizedError} = 401, \texttt{SystemError} = 500. \texttt{statusCode} is \texttt{readonly}.
\end{invariant}

\begin{invariant}
\texttt{isAppError}, \texttt{getErrorMessage}, and \texttt{formatError} are total over \texttt{unknown}: they never throw, regardless of input.
\end{invariant}

\section{Interface}

\begin{lstlisting}[language=TypeScript]
export class AppError extends Error {
  readonly statusCode: number;
  constructor(message: string, statusCode?: number);
}

export class ValidationError extends AppError {
  constructor(message: string); // statusCode = 400
}

export class NotFoundError extends AppError {
  constructor(message: string); // statusCode = 404
}

export class UnauthorizedError extends AppError {
  constructor(message?: string); // statusCode = 401
}

export class SystemError extends AppError {
  constructor(message: string); // statusCode = 500
}

export function isAppError(error: unknown): error is AppError;
export function getErrorMessage(error: unknown): string;
export function formatError(error: unknown): {
  name?: string;
  message: string;
  stack?: string;
};
\end{lstlisting}

\section{Scenarios}

\begin{scenario}{Subclass carries correct status and name}
\Given a new \texttt{ValidationError("bad input")} \\
\When its fields are read \\
\Then \texttt{statusCode === 400} \\
\And \texttt{name === "ValidationError"} \\
\And \texttt{message === "bad input"}.
\end{scenario}

\begin{scenario}{Instanceof hierarchy}
\Given a \texttt{NotFoundError} instance \texttt{e} \\
\When checked with \texttt{instanceof} \\
\Then \texttt{e instanceof NotFoundError} is \texttt{true} \\
\And \texttt{e instanceof AppError} is \texttt{true} \\
\And \texttt{e instanceof Error} is \texttt{true}.
\end{scenario}

\begin{scenario}{isAppError narrows unknown}
\Given an \texttt{unknown} value that is a \texttt{SystemError} instance \\
\When passed to \texttt{isAppError} \\
\Then the guard returns \texttt{true} \\
\And TypeScript narrows the value to \texttt{AppError} in the true branch.
\end{scenario}

\begin{scenario}{getErrorMessage on arbitrary input}
\Given the inputs \texttt{new Error("x")}, \texttt{"y"}, \texttt{42}, \texttt{null} \\
\When each is passed to \texttt{getErrorMessage} \\
\Then it returns \texttt{"x"}, \texttt{"y"}, \texttt{"42"}, \texttt{"null"} respectively \\
\And the function never throws.
\end{scenario}

\begin{scenario}{formatError shape}
\Given an \texttt{AppError("oops")} \\
\When passed to \texttt{formatError} \\
\Then the returned object has \texttt{name}, \texttt{message}, and \texttt{stack} string fields \\
\And \texttt{name === "AppError"}.
\end{scenario}

\section{Acceptance Criteria}

\begin{ac}{AC-1}
\texttt{new AppError("m").statusCode === 500} and \texttt{new AppError("m", 418).statusCode === 418}.
\end{ac}

\begin{ac}{AC-2}
Each subclass sets the expected default \texttt{statusCode}: \texttt{ValidationError} = 400, \texttt{NotFoundError} = 404, \texttt{UnauthorizedError} = 401, \texttt{SystemError} = 500.
\end{ac}

\begin{ac}{AC-3}
For every subclass \texttt{C}, \texttt{(new C(...)).name === "C"} (exact class-name string match).
\end{ac}

\begin{ac}{AC-4}
Every subclass instance satisfies \texttt{instance instanceof AppError \&\& instance instanceof Error}.
\end{ac}

\begin{ac}{AC-5}
\texttt{isAppError(x)} returns \texttt{true} iff \texttt{x} is an instance of \texttt{AppError} or a subclass; otherwise \texttt{false}. It never throws on any \texttt{unknown} input.
\end{ac}

\begin{ac}{AC-6}
\texttt{getErrorMessage(x)} returns \texttt{x.message} when \texttt{x instanceof Error}, otherwise \texttt{String(x)}, for all tested inputs including \texttt{undefined}, \texttt{null}, primitives, and objects.
\end{ac}

\begin{ac}{AC-7}
\texttt{formatError(x)} returns an object with a \texttt{message: string} field for all inputs; when \texttt{x instanceof Error}, the object additionally contains \texttt{name} and \texttt{stack} matching \texttt{x}.
\end{ac}

\begin{ac}{AC-8}
\texttt{UnauthorizedError} instantiated with no argument has \texttt{message === "Unauthorized"} and \texttt{statusCode === 401}.
\end{ac}

\section{Traceability}

\begin{center}
\begin{tabular}{@{}lll@{}}
\toprule
AC & Test file & Test name \\
\midrule
AC-1 & \texttt{packages/shared/tests/app-error.spec.ts} & \texttt{AppError default and override statusCode} \\
AC-2 & \texttt{packages/shared/tests/app-error.spec.ts} & \texttt{subclass default statusCodes} \\
AC-3 & \texttt{packages/shared/tests/app-error.spec.ts} & \texttt{name equals class name} \\
AC-4 & \texttt{packages/shared/tests/app-error.spec.ts} & \texttt{instanceof hierarchy holds} \\
AC-5 & \texttt{packages/shared/tests/app-error.spec.ts} & \texttt{isAppError narrows unknown} \\
AC-6 & \texttt{packages/shared/tests/app-error.spec.ts} & \texttt{getErrorMessage total over unknown} \\
AC-7 & \texttt{packages/shared/tests/app-error.spec.ts} & \texttt{formatError shape for Error and non-Error} \\
AC-8 & \texttt{packages/shared/tests/app-error.spec.ts} & \texttt{UnauthorizedError default message} \\
\bottomrule
\end{tabular}
\end{center}

\section{Open Questions}
\begin{itemize}[leftmargin=1.2em]
  \item Should \texttt{AppError} carry an optional \texttt{cause} field (ES2022) for error chaining, and should \texttt{formatError} surface it?
  \item Should a \texttt{ConflictError} (409) and \texttt{RateLimitError} (429) be added, or kept out to preserve the small, closed hierarchy?
  \item Should \texttt{formatError} redact \texttt{stack} in production environments, or is that the logger's responsibility?
\end{itemize}

\end{document}
